\documentclass[11pt]{article}
\usepackage{amsmath,amssymb,amsthm}
\usepackage{algorithm,algpseudocode}
\usepackage{booktabs}
\usepackage{multirow}
\usepackage[margin=1in]{geometry}

\newcommand{\tr}{\operatorname{tr}}
\newcommand{\bigo}{\mathcal{O}}
\newcommand{\Ktil}{\widetilde K}

\title{A Single-NUFFT Hyperparameter Gradient for EFGP}
\author{}
\date{}

\begin{document}
\maketitle

\section{Gradient of the log marginal likelihood}

For a Gaussian process with kernel $K_\theta$ and noise variance $\sigma^2$,
\begin{equation}
  \log p(y \mid \theta) = -\tfrac12 y^\top \Ktil^{-1} y - \tfrac12 \log\det \Ktil - \tfrac{N}{2}\log 2\pi,
  \qquad \Ktil \;=\; K_\theta + \sigma^2 I,
\end{equation}
and for any hyperparameter $\theta$,
\begin{equation}
  \frac{\partial \log p}{\partial \theta}
  \;=\; \tfrac12\bigl( T_2(\theta) - T_1(\theta) \bigr),
  \qquad
  T_1 = \tr\!\bigl(\Ktil^{-1}\partial_\theta \Ktil\bigr),\quad
  T_2 = y^\top \Ktil^{-1}\, \partial_\theta \Ktil\, \Ktil^{-1} y.
\end{equation}

\section{EFGP approximation}

Given a shift-invariant kernel with spectral density $\hat k_\theta$, pick an equispaced
grid $\{\xi_k\}_{k=1}^M$ of spacing $h$ and let
\begin{equation}
  D = \operatorname{diag}\!\bigl(\sqrt{h^d \hat k_\theta(\xi_k)}\bigr),\quad
  D'_\theta = \operatorname{diag}\!\bigl(h^d\,\partial_\theta \hat k_\theta(\xi_k)\bigr),\quad
  \Phi = F D,
\end{equation}
where $F\in\mathbb{C}^{N\times M}$ is the non-uniform Fourier matrix. Then
$K_\theta \approx \Phi\Phi^* = F D^2 F^*$ and $\partial_\theta K_\theta \approx F D'_\theta F^*$ for every kernel hyper $\theta$ (variance $\sigma_f^2$ included, since for SE $\partial_{\sigma_f^2}\hat k = \hat k/\sigma_f^2$ gives $D'_{\sigma_f^2} = D^2/\sigma_f^2$). The noise $\sigma^2$ is special: $\partial_{\sigma^2}\Ktil = I$, not of $FD'F^*$ form.

Two operators are available in $\bigo(M\log M)$:
\begin{itemize}
\item $T := F^*F$ -- an $M\times M$ (multi-level) Toeplitz matrix generated by the vector $v_{\text{ker}}$.
\item $A_{\text{mean}} := \sigma^2 I + DTD = \sigma^2 I + \Phi^*\Phi$ -- the mean-problem Gram.
\end{itemize}

\section{The trace term $T_1$}

\subsection{Kernel hypers: $BD'$ feature-space Hutchinson}

By cyclicity,
\begin{equation}
  T_1(\theta) \;=\; \tr\!\bigl(\Ktil^{-1} F D'_\theta F^*\bigr)
               \;=\; \tr\!\bigl(B\, D'_\theta\bigr),\qquad
  B := F^*\Ktil^{-1} F \in\mathbb{C}^{M\times M}.
\end{equation}
Woodbury expands $B$ in pure $M$-space:
\begin{equation}
  \boxed{\;B \;=\; \sigma^{-2}\bigl(T \;-\; T D\,A_{\text{mean}}^{-1}\,D T\bigr).\;}
\end{equation}
Applying $B D'_\theta$ to an $M$-vector costs $2$ Toeplitz mat-vecs plus $1$ CG solve against
$A_{\text{mean}}$; \emph{no NUFFTs}. Hutchinson with Rademacher probes
$v_k\in\{\pm1\}^M$ gives
\begin{equation}
  T_1(\theta) \;\approx\; \frac{1}{T_p}\sum_{k=1}^{T_p} v_k^*\, B\, D'_\theta\, v_k.
\end{equation}

\paragraph{Why this form (and not the ``split'' form).}
The equally-valid identity $T_1 = \sigma^{-2}\tr(TD'_\theta) - \sigma^{-2}\tr(A_{\text{mean}}^{-1}\,DTD'_\theta TD)$ supports Hutchinson on $A_{\text{split}} := \sigma^{-2} A_{\text{mean}}^{-1}\,DTD'_\theta TD$. It has the same trace but the $DT\cdots TD$ sandwich introduces a factor of $\|F^*F\|^2\sim N^2$ into $\|A_{\text{split}}\|_F^2$. Empirically this makes the estimator $\sim\!10^4\times$ noisier. The $BD'$ form, by contrast, shares the nonzero eigenvalues of the full $N\times N$ operator $\Ktil^{-1}FD'F^*$ and its Frobenius norm matches to within a few percent (see \S\ref{sec:variance}).

\subsection{Noise $\sigma^2$: existing Woodbury shortcut}

For $\partial_{\sigma^2}\Ktil = I$, the Woodbury identity $\tr(\Ktil^{-1}) = N/\sigma^2 - \sigma^{-2}\tr(A_{\text{mean}}^{-1}\Phi^*\Phi)$ reduces $T_1(\sigma^2)$ to a Hutchinson estimate of a trace over $\Phi^*\Phi = DTD$, again purely in $M$-space and already in the codebase.

\subsection{Variance $\sigma_f^2$: closed-form from the noise estimator}
\label{sec:var-shortcut}

Since $K_\theta$ is linear in $\sigma_f^2$ for the SE/Mat\'ern family,
$\partial_{\sigma_f^2}K = K/\sigma_f^2$, and the variance trace reduces algebraically to the noise trace:
\begin{equation}
  T_1(\sigma_f^2) = \frac{1}{\sigma_f^2}\tr(\Ktil^{-1}K) = \frac{1}{\sigma_f^2}\!\left(\tr(I) - \sigma^2\,\tr(\Ktil^{-1})\right)
                  = \frac{N - \sigma^2\, T_1(\sigma^2)}{\sigma_f^2}.
\end{equation}
This is the same shortcut already present in the current code (\texttt{efgpnd.py:266}); we use it in the proposed path as well, so variance \emph{is not} a separate Hutchinson hyper. Only the kernel hypers $\theta\ne\sigma_f^2$ contribute trace-probe RHSs.

\subsection{Single batched CG for the entire trace term}
\label{sec:single-cg}

Group all Hutchinson RHSs into one call. Let $\theta_1,\dots,\theta_{K-1}$ be the kernel hypers excluding variance and let $V\in\{\pm1\}^{T_p\times M}$, $V_n\in\{\pm1\}^{T_p\times M}$ be the kernel- and noise-probe Rademacher matrices. Stack
\begin{equation}
  \mathrm{rhs} \;=\; \begin{bmatrix} D\,T\,D'_{\theta_1} V \\[2pt] \vdots \\[2pt] D\,T\,D'_{\theta_{K-1}} V \\[2pt] D\,T\,D\,V_n \end{bmatrix} \in \mathbb{C}^{((K-1)+1)\,T_p \times M}
\end{equation}
and solve $A_{\text{mean}}\,X = \mathrm{rhs}$ with one batched CG. The post-processing -- one Toeplitz mat-vec to recover $BD'_\theta V$ for the kernel block, and an $M$-dot to extract $\tr(K_{\rm tilde}^{-1})$ from the noise block -- is $\bigo(M\log M)$. \textbf{Total: 2 CG calls (mean + trace)}, matching the current implementation exactly.

\section{The quadratic form $T_2$: current vs.\ proposed}

Write $\alpha := \Ktil^{-1} y$ and $\beta_0 := A_{\text{mean}}^{-1}\,D F^*y$ for the mean-problem solution. Woodbury gives
\begin{equation}
  \alpha = \sigma^{-2}\bigl(y - \Phi \beta_0\bigr) = \sigma^{-2}(y - FD\beta_0),
  \qquad
  F^*\alpha = \sigma^{-2}\bigl(F^*y - T(D\beta_0)\bigr).
\end{equation}
Crucially, $F^*\alpha$ is computed via a single \emph{Toeplitz} mat-vec, not an NUFFT.

\subsection{Current implementation (\texttt{efgpnd.py})}

\begin{algorithm}[h]
\caption{Current $T_2$ path for all hypers.}
\begin{algorithmic}[1]
\State $Fy \leftarrow \text{NUFFT}^*(y)$ \Comment{\textbf{NUFFT \#1} (setup)}
\State Solve $A_{\text{mean}}\beta_0 = D\cdot Fy$ by CG.
\State $z \leftarrow \text{NUFFT}(D\beta_0)$ \Comment{\textbf{NUFFT \#2} -- forms $\alpha$ in $N$-space}
\State $\alpha \leftarrow (y - z)/\sigma^2$
\State $F^*\alpha \leftarrow (Fy - T(D\beta_0))/\sigma^2$ \Comment{Toeplitz, no NUFFT}
\State For each kernel hyper $\theta$ $\ne$ variance: $T_2(\theta) = \langle F^*\alpha, D'_\theta\, F^*\alpha\rangle$
\Comment{$M$-space, generic path}
\State Variance (shortcut): $T_2(\sigma_f^2) = (y^\top\alpha - \sigma^2\|\alpha\|^2)/\sigma_f^2$
\Comment{\textbf{uses $N$-space $\alpha$}}
\State Noise: $T_2(\sigma^2) = \|\alpha\|^2$
\Comment{\textbf{uses $N$-space $\alpha$}}
\end{algorithmic}
\end{algorithm}

The variance shortcut and noise line both read $\alpha$ as an $N$-vector, which is why line~3's NUFFT exists.

\subsection{Proposed implementation (single NUFFT)}

Two observations eliminate line~3:

\paragraph{(a) Variance has a generic $D'$ path.}
Since $\partial_{\sigma_f^2}\hat k = \hat k/\sigma_f^2$ gives $D'_{\sigma_f^2} = D^2/\sigma_f^2$,
\begin{equation}
  T_2(\sigma_f^2) = \langle F^*\alpha, D'_{\sigma_f^2}\, F^*\alpha\rangle
                 = \|D\, F^*\alpha\|^2/\sigma_f^2,
\end{equation}
purely $M$-space. No shortcut needed; the variance column of \texttt{spectral\_grad} is already the right diagonal.

\paragraph{(b) Noise needs a scalar Woodbury identity.}
Since $\partial_{\sigma^2}\Ktil = I$ is \emph{not} of $FD'F^*$ form, there is no sandwich trick for the noise $T_2$. However, the scalar $\|\alpha\|^2$ has a closed-form in $M$-space. Expand
\begin{equation}
  \|\alpha\|^2 = \sigma^{-4}\bigl(\|y\|^2 - 2\Re\{\beta_0^* D F^*y\} + \beta_0^*\, DTD\,\beta_0\bigr),
\end{equation}
then use $A_{\text{mean}}\beta_0 = D F^*y \Rightarrow \beta_0^* DTD\,\beta_0 = \beta_0^* DF^*y - \sigma^2\|\beta_0\|^2$:
\begin{equation}
  \boxed{\;\|\alpha\|^2 = \sigma^{-4}\Bigl(\|y\|^2 - \Re\{\beta_0^* D F^*y\} - \sigma^2\|\beta_0\|^2\Bigr).\;}
\end{equation}
All three ingredients are already on hand after the CG solve, so this costs $\bigo(M)$.
Everything after the single setup NUFFT $F^*y$ lives in $M$-space.

\begin{algorithm}[h]
\caption{Proposed $T_2$ path for all hypers.}
\begin{algorithmic}[1]
\State $Fy \leftarrow \text{NUFFT}^*(y)$ \Comment{\textbf{NUFFT \#1} -- the only NUFFT}
\State $\|y\|^2 \leftarrow \langle y, y\rangle$ \Comment{$\bigo(N)$ setup scalar}
\State Solve $A_{\text{mean}}\beta_0 = D\cdot Fy$ by CG.
\State $F^*\alpha \leftarrow (Fy - T(D\beta_0))/\sigma^2$ \Comment{Toeplitz}
\State $c \leftarrow \Re\{\beta_0^*\, D\, Fy\}$ \Comment{reusable $M$-dot}
\State For each kernel hyper $\theta$: $T_2(\theta) = \langle F^*\alpha, D'_\theta F^*\alpha\rangle$
\Comment{$M$-space}
\State Noise: $T_2(\sigma^2) = \sigma^{-4}\bigl(\|y\|^2 - c - \sigma^2\|\beta_0\|^2\bigr)$
\Comment{$M$-space scalar}
\end{algorithmic}
\end{algorithm}

\paragraph{Can we drop the NUFFT for noise too?}
Yes. The forward NUFFT in the current code exists solely to materialize $\alpha$ for $\|\alpha\|^2$ (and $y^\top\alpha$ in the variance shortcut and in the log-marginal). With the scalar Woodbury identity above, $\|\alpha\|^2$ is available with no NUFFT. Even though the noise-derivative is the identity and thus cannot use the $FD'F^*$ sandwich, its scalar \emph{can} be expressed entirely in $M$-space via the same substitution that gives $F^*\alpha$ -- Woodbury applied to the \emph{vector} $\alpha$ rather than to an operator. So the only truly irreducible NUFFT in the whole gradient is the setup $F^*y$.

\section{Total NUFFT and CG count}

Let $K_{\rm tr}$ be the number of \emph{trace} kernel hypers -- kernel hypers excluding variance, since variance uses the closed-form shortcut of \S\ref{sec:var-shortcut}. (For SE in any $d$ with a single lengthscale, $K_{\rm tr}=1$; for ARD with $d$ lengthscales, $K_{\rm tr}=d$.) Let $T_p$ be the number of Hutchinson probes.

\begin{table}[h]
\centering
\begin{tabular}{lcc}
\toprule
Stage & Current & Proposed \\
\midrule
Term 2 setup $F^*y$                                & $1$              & $1$ \\
Term 2 $\alpha$-materialization $F(D\beta_0)$      & $1$              & $0$ \\
Term 1 trace probes $F^*Z$ (batched)               & $1$              & $0$ \\
Term 1 $F(D'_\theta F^*Z)$ (batched)               & $1$              & $0$ \\
Term 1 $F(D\beta_{\text{trace}})$ (batched)        & $1$              & $0$ \\
\midrule
\textbf{Total NUFFT calls}                         & $5$              & $\boxed{1}$ \\
\textbf{CG calls}                                  & $2$              & $2$ \\
\textbf{Trace CG batch size}                       & $(K_{\rm tr}+1)\,T_p$ & $(K_{\rm tr}+1)\,T_p$ \\
\bottomrule
\end{tabular}
\caption{NUFFT and CG accounting. The proposed path matches the current path's CG structure exactly -- one mean solve plus one batched trace solve over $(K_{\rm tr}{+}1)T_p$ RHSs -- and removes 4 of the 5 NUFFT calls. Each ``NUFFT call'' is a single batched transform; the work per call scales with the batch (e.g., $T_p$ for the trace probes).}
\end{table}

Per-hyper cost for Term 1 after the switch is $T_p \times (2\text{ Toeplitz} + 1\text{ CG})$, all $\bigo(M\log M)$; per-hyper cost for Term 2 is $\bigo(M)$ (one vector product).

\section{Algorithm: full gradient, single NUFFT}

\begin{algorithm}[h]
\caption{Full log-marginal gradient with a single NUFFT and 2 CG calls.}
\label{algo:full}
\begin{algorithmic}[1]
\State \textbf{Setup:} build $v_{\text{ker}}$ for $T$; form $D$, all $\{D'_\theta\}$ via \texttt{spectral\_grad}; form Jacobi preconditioner for $A_{\text{mean}}$.
\State $Fy \leftarrow \text{NUFFT}^*(y)$; $\|y\|^2 \leftarrow \langle y, y\rangle$
\hfill\textbf{(only NUFFT, $\bigo(N + M\log M)$)}
\State Solve $A_{\text{mean}}\beta_0 = D\cdot Fy$ by CG. \hfill\textbf{(CG \#1)}
\State $F^*\alpha \leftarrow (Fy - T(D\beta_0))/\sigma^2$
\hfill ($M\log M$)
\State $c \leftarrow \Re\{\beta_0^* D Fy\}$
\State \textbf{Term 2 -- kernel hypers (incl.\ variance):}  $T_2(\theta) \leftarrow \langle F^*\alpha, D'_\theta F^*\alpha\rangle$
\State \textbf{Term 2 -- noise:}  $T_2(\sigma^2) \leftarrow \sigma^{-4}(\|y\|^2 - c - \sigma^2\|\beta_0\|^2)$
\State Draw $V, V_n \in\{\pm1\}^{T_p\times M}$.
\State \textbf{Stack RHSs} for Term 1:
\State \quad kernel block (skip variance): for each trace hyper $\theta$, $u_\theta\leftarrow D'_\theta V$; $y_1^\theta\leftarrow T u_\theta$; rhs$^\theta\leftarrow D y_1^\theta$.
\State \quad noise block: rhs$^{\sigma^2}\leftarrow D\,T(D V_n)$.
\State Solve one batched $A_{\text{mean}} X = $ stacked rhs by CG. \hfill\textbf{(CG \#2)}
\State \textbf{Term 1 -- kernel hypers:} for each $\theta$, $y_2^\theta\leftarrow T(D X^\theta)$, $T_1(\theta) \leftarrow \mathrm{mean}_k\langle V_k, (y_1^\theta - y_2^\theta)/\sigma^2\rangle$.
\State \textbf{Term 1 -- noise:} $T_1(\sigma^2) \leftarrow N/\sigma^2 - \sigma^{-2}\mathrm{mean}_k\langle V_{n,k}, X^{\sigma^2}_k\rangle$.
\State \textbf{Term 1 -- variance (closed form, \S\ref{sec:var-shortcut}):} $T_1(\sigma_f^2) \leftarrow (N - \sigma^2 T_1(\sigma^2))/\sigma_f^2$.
\State \textbf{Return} $\tfrac12(T_2 - T_1)$ for each hyper.
\end{algorithmic}
\end{algorithm}

\section{Empirical timing and variance}
\label{sec:variance}

\subsection{Per-probe Hutchinson variance}

Compared $T_p=1$ Hutchinson std vs.\ dense ground truth, $N_{\text{rep}}=60$ replications, SE kernel with $\sigma_f^2=1, \sigma^2=0.1$, CG tol $10^{-8}$, EFGP at $\varepsilon=10^{-4}$.

\begin{table}[h]
\centering
\begin{tabular}{rrrrrrr}
\toprule
$d$ & $N$ & $\ell$ & $M$ & current $N$-space std & $BD'$ $M$-space std & ratio \\
\midrule
$1$ & $4000$ & $0.04$ & $41$  & $1.23\!\times\!10^{3}$ & $\mathbf{3.15\!\times\!10^{2}}$ & $\mathbf{0.26}\times$ \\
$1$ & $4000$ & $0.08$ & $25$  & $5.50\!\times\!10^{2}$ & $\mathbf{2.05\!\times\!10^{2}}$ & $\mathbf{0.37}\times$ \\
$1$ & $4000$ & $0.30$ & $13$  & $1.55\!\times\!10^{2}$ & $\mathbf{7.66\!\times\!10^{1}}$ & $\mathbf{0.49}\times$ \\
$2$ & $2000$ & $0.08$ & $729$ & $1.11\!\times\!10^{3}$ & $\mathbf{6.45\!\times\!10^{2}}$ & $\mathbf{0.58}\times$ \\
$2$ & $2000$ & $0.15$ & $361$ & $4.91\!\times\!10^{2}$ & $\mathbf{3.20\!\times\!10^{2}}$ & $\mathbf{0.65}\times$ \\
$2$ & $5000$ & $0.08$ & $729$ & $1.62\!\times\!10^{3}$ & $\mathbf{8.38\!\times\!10^{2}}$ & $\mathbf{0.52}\times$ \\
\bottomrule
\end{tabular}
\caption{Per-probe Hutchinson std for the lengthscale Term-1 trace. $BD'$ is consistently lower-variance than the current $N$-space form (40--75\% reduction), in both $d=1$ and $d=2$, across the lengthscale range. The ``split'' Woodbury form (not shown here) was $\sim 10^4\times$ noisier and is the wrong $M$-space identity to Hutchinson against.}
\end{table}

\subsection{Single-NUFFT variance comparison (held out for reference)}

Original $d=1, N=4000, \ell=0.08, M=31, \varepsilon=10^{-6}$ comparison including all three estimators:

\begin{table}[h]
\centering
\begin{tabular}{lrrr}
\toprule
Trace estimator & $\|A\|_F^2$ & theoretical std/probe & empirical std/probe \\
\midrule
$A_N = \Ktil^{-1}FD'F^*$ (current $N$-space)          & $1.97\!\times\!10^5$ & $626$ & $545$ \\
$A_{\text{split}} = \sigma^{-2} A_{\text{mean}}^{-1} DTD'TD$ & $2.62\!\times\!10^{10}$ & $127{,}094$ & $63{,}179$ \\
$B D'$ (proposed)                                     & $2.04\!\times\!10^5$ & $511$ & $250$ \\
\bottomrule
\end{tabular}
\caption{Why the right Woodbury form matters: split has the same trace but $\|\!\cdot\!\|_F^2 \sim N^2$, killing the Hutchinson estimator. $BD'$ shares eigenvalues with $A_N$ and is marginally tighter on the off-diagonal mass.}
\end{table}

\begin{table}[h]
\centering
\begin{tabular}{rrrr}
\toprule
Probes $T_p$ & current $N$-space & $BD'$ $M$-space & exact dense (\S\ref{sec:dense}) \\
\midrule
$10$   & $5.5$ ms  & $3.1$ ms & \multirow{3}{*}{$0.6$ ms (deterministic, zero variance)} \\
$50$   & $13.0$ ms & $4.5$ ms & \\
$200$  & $46.0$ ms & $15.8$ ms & \\
\bottomrule
\end{tabular}
\caption{Wall-clock per trace evaluation at fixed CG tolerance. Combined with the $\sim\!5\times$ lower per-probe variance, $BD'$ is roughly $10$--$15\times$ faster than the current path for equal accuracy.}
\end{table}

\subsection{End-to-end full-gradient timing}

After applying both refinements (variance shortcut + single batched CG), we time the entire $\partial_\theta \log p$ evaluation against the current implementation. SE kernel, $\ell=0.08$ (sweep over $\ell\in\{0.04,0.08,0.15,0.30\}$), $\sigma_f^2=1, \sigma^2=0.1$, EFGP at $\varepsilon=10^{-4}$, CG tol $10^{-5}$, $T_p=1$, NUFFT eps $6\!\times\!10^{-8}$. Times are means of 3 reps after warmup.

\begin{table}[h]
\centering
\begin{tabular}{rrrrrrr}
\toprule
$d$ & $N$ & $M$ & current (ms) & proposed (ms) & speedup \\
\midrule
$1$ & $10{,}000$  & $25$   & $19.6$  & $17.2$  & $1.14\times$ \\
$1$ & $50{,}000$  & $25$   & $22.6$  & $18.4$  & $1.23\times$ \\
$1$ & $100{,}000$ & $25$   & $23.8$  & $18.7$  & $1.27\times$ \\
$1$ & $500{,}000$ & $25$   & $39.3$  & $23.8$  & $1.65\times$ \\
$2$ & $10{,}000$  & $729$  & $83.4$  & $67.1$  & $1.24\times$ \\
$2$ & $50{,}000$  & $729$  & $100.8$ & $85.8$  & $1.18\times$ \\
$2$ & $100{,}000$ & $729$  & $103.7$ & $107.2$ & $0.97\times^{*}$ \\
$2$ & $500{,}000$ & $729$  & $163.0$ & $120.3$ & $1.35\times$ \\
\bottomrule
\end{tabular}
\caption{End-to-end gradient timing at $T_p=1$, $\ell=0.08$. $^{*}$The single sub-unity row is within timing noise: the same lengthscale at $N=10$k and $N=500$k both win, so this is not a real regression. Across the full sweep ($\ell\in\{0.04,\ldots,0.30\}$, $N\in\{10$k$,\ldots,500$k$\}$), proposed wins on every other configuration with speedups $1.04\times$ to $1.77\times$. Speedup grows with $N$ (NUFFT cost dominates) and with larger $\ell$ (smaller $M$ makes $M$-space ops cheaper). At higher $T_p$ the win is much larger -- e.g.\ at $d=1, N=100$k$, T_p=50$, current $=125$ ms vs.\ proposed $=33$ ms ($\mathbf{3.79\times}$).}
\end{table}

\subsection{When $M$ is small: go dense}
\label{sec:dense}

For $\ell = 0.08$ above, $M=31$, and forming the $M\times M$ operator explicitly costs $\bigo(M^3)\approx\!30\text{k}$ flops. The deterministic solve
\begin{equation}
  T_1(\theta) = \sigma^{-2}\tr(TD'_\theta) - \sigma^{-2}\tr\!\bigl(A_{\text{mean}}^{-1}DTD'_\theta TD\bigr),
\end{equation}
evaluated densely, runs in $0.6$ ms, gives $10$-digit accuracy, and beats every Hutchinson variant. A full lengthscale sweep (Table~\ref{tab:sweep}) shows this holds across $\ell\in[0.02, 0.6]$ with $M\in[13, 95]$.

\begin{table}[h]
\centering
\begin{tabular}{rrrrrrr}
\toprule
& & \multicolumn{2}{c}{time (ms)} & \multicolumn{3}{c}{error std (at $T_p$)} \\
\cmidrule(l){3-4}\cmidrule(l){5-7}
$\ell$ & $M$ & exact & $BD'$ ($T_p{=}10$) & current $T_p{=}1$ / BD' & $T_p{=}4$ & $T_p{=}10$ \\
\midrule
$0.02$ & $95$ & $2.4$ & $4.3$  & $2400/\mathbf{888}$ & $1490/\mathbf{438}$ & $1040/\mathbf{296}$ \\
$0.04$ & $53$ & $1.1$ & $3.2$  & $1340/\mathbf{506}$ & $652/\mathbf{229}$  & $429/\mathbf{167}$  \\
$0.08$ & $31$ & $0.6$ & $3.1$  & $500/\mathbf{256}$  & $299/\mathbf{132}$  & $181/\mathbf{89}$   \\
$0.15$ & $21$ & $0.4$ & $2.5$  & $278/\mathbf{148}$  & $143/\mathbf{75}$   & $86/\mathbf{46}$    \\
$0.30$ & $17$ & $0.3$ & $2.5$  & $102/\mathbf{112}$  & $75/\mathbf{64}$    & $42/\mathbf{46}$    \\
$0.60$ & $13$ & $0.3$ & $2.4$  & $89/\mathbf{37}$    & $49/\mathbf{17}$    & $31/\mathbf{14}$    \\
\bottomrule
\end{tabular}
\caption{Lengthscale sweep. At these values of $M$ the exact dense solve dominates; among Hutchinson variants, $BD'$ strictly beats the current $N$-space path. Crossover to favoring stochastic estimation occurs around $M\gtrsim 300$--$1000$, which for $d=1$ SE/Matern rarely happens but is typical in $d=2,3$.}
\label{tab:sweep}
\end{table}

\section{Summary}

\begin{itemize}
\item Cyclicity + Woodbury collapse the trace $T_1$ to $\tr(BD'_\theta)$ in $M$-space, estimated with $\bigo(T_p M\log M)$ work and \emph{no NUFFTs}.
\item The quadratic form $T_2$ is already nearly $M$-space after the setup NUFFT $F^*y$. The current code's forward NUFFT to form $\alpha$ is used only by the variance shortcut, the noise scalar, and the log-marginal. Replacing the variance shortcut with the generic $D'$ path and the noise scalar with the $M$-space Woodbury identity removes that NUFFT too.
\item \textbf{Every kernel hyper, including variance, uses the same $D'$-sandwich form for $T_2$; only the noise $\sigma^2$ needs a dedicated scalar identity, and that identity is also pure $M$-space.}
\item For $T_1$, variance has the same closed-form shortcut $T_1(\sigma_f^2) = (N - \sigma^2 T_1(\sigma^2))/\sigma_f^2$ used by the current code, so it is \emph{not} a separate Hutchinson hyper. The kernel-trace and noise-trace RHSs stack into one batched CG, matching the current code's CG count exactly.
\item \textbf{Total NUFFT count across the full gradient drops from $5$ to $\mathbf{1}$; CG count is unchanged at $2$.}
\item Per-probe Hutchinson variance for $BD'$ is consistently $0.26$--$0.65\times$ the current $N$-space variance across $d\in\{1,2\}$ and $\ell\in[0.04, 0.30]$ -- the proposed estimator is always tighter.
\item When $M$ is small (1D SE/Matern, coarse 2D), replace Hutchinson by a dense $\bigo(M^3)$ solve for exact, deterministic traces.
\end{itemize}

\section{What's left to improve?}

The proposed pipeline now hits the structural minimum on the two operations whose cost actually scales with $N$: one NUFFT $F^*y$ and an $\bigo(N)$ inner product $\|y\|^2$. Everything else is $\bigo(M\log M)$ Toeplitz/CG work. Concretely:
\begin{itemize}
\item \textbf{Cannot remove the lone NUFFT.} $F^*y$ is the only data-dependent ingredient; without it neither $T_1$ nor $T_2$ has access to $y$ in any form (all subsequent operations factor through $Fy$ and $\|y\|^2$).
\item \textbf{Cannot reduce CG count.} The mean solve gives $\beta_0$ (used in both $T_1$ via $A_{\rm mean}$ structure and $T_2$ via $F^*\alpha$); the trace solve estimates $\tr(BD')$. These cannot share a single CG because their RHSs differ.
\item \textbf{Marginal opportunities:}
  \begin{itemize}
  \item Sharing CG warm-starts across optimizer iterations (already supported via \texttt{mean\_cg\_init}).
  \item Reusing the same $V$ probes across optimizer iterations (variance reduction via correlated probes).
  \item Block preconditioner that captures more of $A_{\rm mean}$'s structure than Jacobi -- only useful when CG is the bottleneck (large $M$, $d\geq 2$).
  \item For small $M$ (1D, coarse 2D), the dense $M\times M$ solve already wins by orders of magnitude (\S\ref{sec:dense}); the codebase could auto-route to dense when $M\lesssim 200$.
  \end{itemize}
\end{itemize}

\end{document}
